\begin{abstract}
Malware behavior analysis at scale requires transparent, low-perturbation
observation of software execution. Software tracers such as strace, eBPF, and
QEMU run inside the target environment, making them detectable by evasive
workloads and introducing perturbation that can alter behavior. Hardware-assisted
tracing, such as ARM ETM and Intel PT, provides an alternative observation path
below the OS, but produces raw architectural events that lack OS-level semantics
and must be carefully reconstructed to answer security questions.

We propose \rvmt, a hardware-assisted behavior tracing system for RISC-V Linux
that connects a CVA6/Genesys2 hardware trace path with an offline behavior
reconstruction pipeline. \rvmt captures architectural events---retire, trap,
syscall entry/return, privilege transition, branch/jump, context, and marker
events---below the OS, then joins them with ELF metadata and runtime process maps
while explicitly labeling each reconstructed field with its provenance source. Our
evaluation demonstrates that the trace pipeline passes 50 directed correctness
cases, 10 seeded-random event-sequence cases, and 9 negative-sensitivity cases.
Across 12 controlled workload artifacts, 288/288 provenance-tagged semantic
fields pass source-label validation. The code-attribution pipeline attributes
108/108 key PC events to the target ELF/function level and passes 6/6 mapping
edge cases covering PIE/ASLR, dynamic loading, fork/exec, and stripped binaries.
For a controlled anti-debug-like workload where software tracing is visible or
perturbs \texttt{ptrace}, \rvmt reconstructs the syscall sequence and hardware
\texttt{ARG\_MEM} string needed to identify the behavior. We also report the
runtime and transport evidence boundaries: current measurements include board
UART runtime smoke, FPGA trace-cycle marker windows, and a BRAM-derived target
for future streaming/DMA transport, but not cycle-level slowdown or production
throughput. The artifact package rooted at \artifactroot records 63 summary
artifacts, 240 package files, and a 2,033-file raw-artifact archive with SHA-256
checks for independent verification.
\end{abstract}
